% META: {"id": "TNSI-ARCH-CO-005", "chapitre": "TNSI-ARCHITECTURES-MATERIELLES-SY", "type_objet": "corrige", "exercice_ref": "TNSI-ARCH-EX-005", "capacites": ["T-ARCH-01", "T-ARCH-02A"], "capacites_codes": ["C1", "C2"], "mode_creation": "ex_nihilo", "sources_inspiration": [], "status": "generated"}
% BEGIN-VERIFY
% fiche = {"processeurs": 8, "memoire_vive_Go": 8,
%          "accelerateurs": ["graphique", "reseau de neurones", "signal audio"],
%          "interfaces": ["Wi-Fi", "Bluetooth", "USB", "modem 5G"]}
% assert 1 + 1 + len(fiche["interfaces"]) + len(fiche["accelerateurs"]) == 9
% assert "modem 5G" in fiche["interfaces"]
% def creer(arbre, parent, enfant):
%     arbre = {p: list(f) for p, f in arbre.items()}
%     arbre.setdefault(parent, [])
%     arbre[parent].append(enfant)
%     arbre.setdefault(enfant, [])
%     return arbre
% arbre = {"init": []}
% creations = [("init", "session"), ("session", "terminal"),
%              ("session", "navigateur"), ("navigateur", "onglet1"),
%              ("navigateur", "onglet2")]
% for parent, enfant in creations:
%     arbre = creer(arbre, parent, enfant)
% assert len(arbre) == 6 and len(creations) == 5
% assert len(arbre) == len(creations) + 1
% assert sorted(arbre) == ["init", "navigateur", "onglet1", "onglet2", "session", "terminal"]
% assert arbre["session"] == ["terminal", "navigateur"]
% assert arbre["init"] == ["session"]
% def descendants(arbre, racine):
%     total = 0
%     for enfant in arbre[racine]:
%         total = total + 1 + descendants(arbre, enfant)
%     return total
% assert descendants(arbre, "init") == 5
% assert descendants(arbre, "session") == 4
% assert descendants(arbre, "navigateur") == 2
% assert descendants(arbre, "terminal") == 0
% assert descendants(arbre, "init") == len(arbre) - 1
% END-VERIFY
\begin{corrige}{TNSI-ARCH-EX-005}
\textbf{1.} Le classement est immédiat :

\begin{center}
\begin{tabular}{ll}
\hline
catégorie & éléments de la fiche \\
\hline
processeurs & les $8$ c\oe{}urs \\
mémoire & les $8$~Go de mémoire vive \\
interfaces & Wi-Fi, Bluetooth, USB, modem 5G \\
accélérateurs & graphique, réseau de neurones, signal audio \\
\hline
\end{tabular}
\end{center}

En comptant les processeurs pour un composant et la mémoire pour un autre, cela
fait $1+1+4+3=9$ composants intégrés sur une seule puce. Chacun d'eux occupait
autrefois un circuit distinct sur une carte mère.

\textbf{2.} Les trois avantages sont la \textbf{compacité} (une seule puce au
lieu d'une carte), la \textbf{consommation énergétique réduite} (les distances
entre composants sont très courtes, donc les pertes liées à leurs échanges le
sont aussi) et le \textbf{coût de fabrication} (une puce produite en grande série
revient moins cher qu'un assemblage).

C'est le deuxième qui explique l'autonomie : sur une carte mère, faire circuler
un signal du processeur vers la mémoire coûte de l'énergie proportionnellement à
la distance parcourue, et ces échanges sont incessants. Les rapprocher de
quelques millimètres change l'ordre de grandeur de la facture énergétique.

\textbf{3.} Il faut remplacer \textbf{toute la puce}, donc en pratique changer
le téléphone ou sa carte principale — avec les huit c\oe{}urs, les huit gigaoctets
de mémoire et les trois accélérateurs qui fonctionnaient parfaitement. Sur
l'ordinateur de bureau, on remplace la seule carte réseau, pour une fraction du
prix.

La propriété perdue est la \textbf{modularité} : la capacité à remplacer ou faire
évoluer un composant indépendamment des autres. C'est le contrepoids exact des
trois avantages, et il n'est pas mince — il pèse sur la durée de vie des
appareils.

\textbf{4.} Le dictionnaire associe à chaque processus ses enfants :
\begin{python}
def creer(arbre, parent, enfant):
    arbre = {p: list(f) for p, f in arbre.items()}
    arbre.setdefault(parent, [])
    arbre[parent].append(enfant)
    arbre.setdefault(enfant, [])
    return arbre
\end{python}
En partant de \lstinline!{"init": []}! et en appliquant les cinq créations, on
obtient
\lstinline!{"init": ["session"], "session": ["terminal", "navigateur"],!
\lstinline!"terminal": [], "navigateur": ["onglet1", "onglet2"],!
\lstinline!"onglet1": [], "onglet2": []}!.

\textbf{5.} Il y a $6$ processus pour $5$ créations : le nombre de processus
vaut toujours le nombre de créations plus un.

La raison est que chaque création ajoute \emph{exactement} un processus, et que
le premier — \lstinline{init} — n'a été créé par personne. Tout processus autre
que lui a un parent et un seul, ce qui fait de l'ensemble un arbre : $n$ sommets,
$n-1$ arêtes. C'est aussi ce qui garantit qu'aucun processus n'est orphelin de
l'arborescence.

\textbf{6.} La fonction descend récursivement :
\begin{python}
def descendants(arbre, racine):
    total = 0
    for enfant in arbre[racine]:
        total = total + 1 + descendants(arbre, enfant)
    return total
\end{python}
On obtient $5$ pour \lstinline{init} — soit tous les autres processus, ce qui
confirme la relation de la question précédente —, $2$ pour
\lstinline{navigateur} et $0$ pour \lstinline{terminal}, qui n'a rien créé.
Fermer le navigateur reviendrait à terminer trois processus, et non un seul.
\end{corrige}
